\begin{problem}[理论考试][中性粒子分析器]{29}
  中性粒子分析器(Neutral-Particle Analyser)是核聚变研究中测量快离子温度及其能量分布的重要设备。其基本原理如图所示，通过对高能量（200 eV ~ 30 KeV）中性原子（它们容易穿透探测区中的电磁区域）的能量和动量的测量，可诊断曾与这些中性原子充分碰撞过的离子的性质。

  为了测量中性原子的能量分布，首先让中性原子电离，然后让离子束以 $\theta$ 角入射到间距为 $d$、电压为 $V$ 的平行板电极组成的区域，经电场偏转后离开电场区域，在保证所测量离子不碰到上极板的前提下，通过测量入射孔A和出射孔B间平行于极板方向的距离 $l$ 来决定离子的能量。设A与下极板的距离为 $h_1$，B与下极板的距离为 $h_2$，已知离子所带电荷为 $q$：
  \begin{subproblem}
    推导离子能量 E 与 l 的关系，并给出离子在极板内垂直于极板方向的最大飞行距离。
  \end{subproblem}
  \begin{subproblem}
    被测离子束一般具有发散角 $\Delta \alpha (\Delta \alpha = \theta)$。为了提高测量的精度，要求具有相同能量 $E$，但入射方向在 $\Delta \alpha$ 范围内变化的离子在同一小孔 B 处射出，求 $h_2$ 的表达式；并给出此时能量 $E$ 与 $l$ 的关系。
  \end{subproblem}
  \begin{subproblem}
    为了提高离子能量的分辨率，要求具有量程上限能量的离子刚好落在设备允许的 $l$ 的最大值 $l_{\mathrm{max}}$ 处，同时为了减小设备的体积，在满足测量要求的基础上，要求极板间距 $d$ 尽可能小，利用上述第（2）问的结果，求 $d$ 的表达式；若 $\theta = 30^{\circ}$，结果如何？
  \end{subproblem}
  \begin{subproblem}
    为了区分这些离子的质量，请设计后续装置，给出相应的原理图和离子质量表达式。
  \end{subproblem}
\end{problem}